% AI-produced by GPT-6 Astra (OpenAI, ChatGPT Pro mode). No priority claim.
\subsection{A compact inverse-limit reformulation of a stronger Ramsey condition}

This subsection isolates a possible strengthening of totality. It does
not establish that strengthening. A compact inverse-limit argument
cannot be used to prove the main implication without supplying the
missing hypothesis identified below.

\begin{definition}[Proved: agreement of the two Ramsey formulations for finite hom-sets]\label{CD:FPR}
A small category $A$ has the finite-pattern morphism Ramsey property if,
for every $p,q\in A$, finite subset $S\subset A(p,q)$, and positive
integer $c$, there is $r\in A$ such that every coloring
\[
 \chi:A(p,r)\longrightarrow\{1,\ldots,c\}
\]
admits $g:q\to r$ for which $\chi$ is constant on $gS$.
When all hom-sets are finite, this is equivalent to the morphism Ramsey
property in Theorem~\ref{OG:Ramsey-exact}: take $S=A(p,q)$.
\end{definition}

\begin{theorem}[Proved: compact inverse limits and finite-pattern Ramsey]
\label{CD:compact-Ramsey}
Let $A$ be a small nonempty category with common successors. The following
are equivalent.
\begin{enumerate}
\item Every contravariant diagram $X:A^{\mathrm{op}}\to
\mathrm{CompHaus}$ of nonempty compact Hausdorff spaces has a nonempty
strict inverse limit.
\item $A$ has the finite-pattern morphism Ramsey property.
\end{enumerate}
No bound on the hom-sets is assumed.
\end{theorem}
\begin{proof}
Suppose (2) fails for $p,q,S,c$. Necessarily $S$ is nonempty. For
each $a$, let $X(a)$ consist of the colorings of $A(p,a)$ that are
nonconstant on $gS$ for every $g:q\to a$. This is a closed subset of
the compact product $\{1,\ldots,c\}^{A(p,a)}$: each forbidden
monochromaticity event concerns finitely many coordinates and is clopen.
The failure of (2) says that every $X(a)$ is nonempty. Precomposition
with $A(p,u)$, for $u:a\to b$, defines a continuous map
$X(b)\to X(a)$, since a monochromatic $gS$ after restriction would be
a monochromatic $(ug)S$ before restriction. Thus $X$ is a diagram
as in (1). A compatible family $(\chi_a)_a$ would satisfy
\[
 \chi_a(h)=\chi_p(1_p)\qquad(h:p\to a).
\]
In particular $\chi_q$ would be constant on $S$, contrary to its
defining condition for $g=1_q$. This proves $(1)\Rightarrow(2)$.

For the converse, first observe that a Ramsey target has a finite
witness set. If $r$ works for $p,q,S,c$, the clopen sets
\[
 U_g=\{\chi:\chi\text{ is constant on }gS\},\qquad g:q\to r,
\]
cover the compact coloring space. A finite subcover gives a finite
$G\subset A(q,r)$ such that some member of $G$ witnesses every
coloring.

It follows that finitely many finite-pattern requirements can be
satisfied simultaneously. To spell this out, let
$S_j\subset A(p_j,q)$ and $c_j$ be given for $1\leq j\leq m$.
Put $r_0=q$. At stage $j$ apply (2) to the finite pattern
\[
 S_j^{(j-1)}=
 \bigcup_{h_i\in G_i,\,1\leq i<j}h_{j-1}\cdots h_1S_j
 \subset A(p_j,r_{j-1}),
\]
where $S_1^{(0)}=S_1$. Obtain a target $r_j$ and a finite witness
set $G_j\subset A(r_{j-1},r_j)$. Given colorings at $r_m$, choose
$h_m\in G_m$ first, then $h_{m-1}\in G_{m-1}$ for the appropriate
pullback coloring, and continue backwards. The composite
$g=h_m\cdots h_1:q\to r_m$ makes every $gS_j$ monochromatic.
Indeed the prefix $h_{j-1}\cdots h_1S_j$ is contained in
$S_j^{(j-1)}$, so the choice at stage $j$ controls every possible
subsequent choice of that prefix.

Now let $X$ be as in (1). Its limit is cut out in the compact product
$\prod_aX(a)$ by the closed equations
\[
 x_a=X(f)x_b\qquad(f:a\to b).
\]
Fix finitely many such equations, indexed by a finite set $E$ of
arrows. For each source $a$ in $E$, fix a neighborhood $U_a$ of the
diagonal in $X(a)\times X(a)$. Choose a finite cover of $X(a)$ by
sets whose Cartesian squares lie in $U_a$, and assign each point to
one containing cover member. Such a cover exists by the compatible
uniformity of a compact Hausdorff space. This gives a finite coloring
of $X(a)$; continuity of this coloring is not required.

Choose a common successor $q$ of the finitely many objects involved
and arrows $j_a:a\to q$. For each source $a$ use the finite pattern
\[
 S_a=\{j_a\}\cup\{j_bf:(f:a\to b)\in E\}.
\]
Apply the simultaneous Ramsey assertion with the cover color counts
to obtain $r$. Choose $z\in X(r)$ and color $h:a\to r$ by the cover
color of $X(h)z$. The resulting $g:q\to r$ makes the assignments
$x_a=X(gj_a)z$ satisfy
\[
 (x_a,X(f)x_b)\in U_a\qquad((f:a\to b)\in E).
\]
Choose arbitrary values for the other coordinates. Thus every finite
set of equations can be satisfied to arbitrary uniform accuracy.
Use closed entourages for each equation: their inverse images in the
product have the finite intersection property, while the intersection
over all such entourages is exactly the equality condition. Compactness
therefore supplies a point satisfying every equation. This proves (1).
\end{proof}

\begin{corollary}[Proved: the compact assertion implies integral exactness]
\label{CD:compact-exact}
A nonempty small category with common successors satisfying either
condition of Theorem~\ref{CD:compact-Ramsey} has exact abelian colimits.
\end{corollary}
\begin{proof}
The proof of Theorem~\ref{OG:Ramsey-exact} uses monochromaticity only
on the specified finite family $\{a_0,\ldots,a_k\}$, so its Ramsey
input can be replaced by Definition~\ref{CD:FPR}. For arbitrary
hom-sets, replace its finite-rank Smith normal form step by extension
of a character into $\mathbb R/\mathbb Z$. For completeness, that
extension needs only divisibility: choose a maximal partial extension
on a subgroup by Zorn's lemma. If its domain is not the whole group,
adjoin an element $x$. If a least positive $n$ has $nx$ in the old
domain, choose an $n$th root of its assigned value; if no such $n$
exists, choose any value for $x$. This extends the character in either
case, a contradiction. The circle argument then gives an integral
affine equalizer for every finite parallel family.

To see directly that these equalizers suffice, two parallel affine
arrows have finite supports with augmentation one. Equalizing all
ordinary arrows in the union of those supports by a further affine
arrow makes the two affine composites equal. Common successors and
this property make $\Aff_\Z(A)$ filtered. The exactness criterion of
\citet[Theorem~3.1 and Proposition~3.3]{IsbellMitchell1976} applies.
\end{proof}

\begin{proposition}[Proved: uniformly bounded finite fibers]
\label{CD:bounded-fibres}
Let $A$ be nonempty, connected, and quasi-filtered. Every contravariant
diagram $X:A^{\mathrm{op}}\to\Set$ of nonempty finite sets with
$|X(a)|\leq N$ for one finite integer $N$ has a nonempty inverse limit.
\end{proposition}
\begin{proof}
Let $A_X$ have the objects of $A$, with parallel arrows identified
precisely when they induce the same function $X(b)\to X(a)$.
Functoriality makes this an equivalence relation compatible with
composition. Conditions I and IV descend from $A$ by choosing
representatives, applying their required square or affine equality
in $A$, and then passing to the quotient. Hence $A_X$ is quasi-filtered.
For each fixed $a$ one has
\[
 |A_X(a,b)|\leq |X(a)|^{|X(b)|}\leq |X(a)|^N
 \qquad\text{for every }b.
\]
By \citet[Proposition~4.5]{IsbellMitchell1976}, every finite parallel
family out of $a$ is equalized by an ordinary arrow in $A_X$.
Connectedness and I give common successors, so $A_X$ is filtered.

The diagram $X$ factors through $A_X^{\mathrm{op}}$, with the same
strict limit. Every finite set of compatibility equations for this
diagram has a solution: filteredness gives a cocone on the finite
diagram in question, and any point at its vertex supplies that
solution. The compatibility conditions are closed in the compact
product of the finite discrete values. They have the finite
intersection property, so their intersection is nonempty.
\end{proof}

\begin{example}[Proved: coefficient pinning can destroy a colimit relation]
\label{CD:pinning-counterexample}
Let $A=\Delta_{\mathrm{inj}}$, which is exact and total, and let
$B(n)$ be the free Boolean algebra on the generators
$x_0,\ldots,x_{n-1}$, also regarded as its associated Boolean ring
with addition given by symmetric difference. An increasing injection $f:n\to m$ acts by
$x_i\mapsto x_{f(i)}$. These are injective unital Boolean-algebra
homomorphisms: the corresponding restriction maps on the sets of
Boolean valuations are surjective.

In $B(2)$ the elements $x_0,x_1$ have the same class in
$\colim_A^{\Set}UB$, since the two maps $1\to2$ carry the sole
generator to them. Nevertheless, with $c=x_0$, the pairs
$(c,x_0)$ and $(c,x_1)$ lie in different components of
$\El_A(UB\times UB)$. Equality of the two coordinates is preserved
and reflected by every structure map; it holds for the first pair
and fails for the second.

The failure persists for the abelian colimit after pinning $c$.
Consider the component $E$ of $\El_A UB$ containing $(2,x_0)$.
Its objects are exactly the pairs $(n,x_j)$. Indeed renaming
variables sends a Boolean expression to a generator only if that
expression was itself a generator, as can be checked on valuations.
On $E$, evaluation at the valuation which sends $x_j$ to $1$ and all
other generators to $0$ defines a natural homomorphism of additive
groups
\[
 B(n)\longrightarrow\F_2.
\]
For an arrow $f:(n,x_j)\to(m,x_{f(j)})$, the target valuation
restricts to the specified source valuation, proving naturality.
At $(2,x_0)$ this map takes $x_0$ to $1$ and $x_1$ to $0$.
Consequently $x_0$ and $x_1$ remain distinct in the abelian colimit
over $E$.

Thus moving to a category of elements in order to hold a multiplicative
coefficient fixed need not preserve even a set-colimit relation
available before that move. In a proposed Boolean-algebra argument,
one must prove the needed lifted tuple relations; totality alone
cannot be cited for their existence. This example does not rule out
a different homological argument exploiting the full Boolean identity.
\end{example}

\begin{obstruction}[Open in this incomplete record: the missing compact implication]
\label{CD:unresolved}
We have not proved that totality implies either condition of
Theorem~\ref{CD:compact-Ramsey}. Nor have we proved or refuted that
exactness implies those conditions. Proposition~\ref{CD:bounded-fibres}
only shows that a finite-valued obstruction for a quasi-filtered
category must have unbounded fiber cardinalities. Signed measures
obtained from affine equalizers cannot be replaced by genuine points
without an additional argument.
\end{obstruction}
